\section{Decisiones y limitaciones}
\label{sec:decisiones}

\begin{table}[H]
\centering\small
\begin{tabularx}{\textwidth}{@{}L L@{}}
\toprule
\textbf{Decisión} & \textbf{Por qué} \\
\midrule
Órdenes discretas en lugar de joystick & El guardia supervisa una ronda, no maneja un vehículo. Un joystick invita a manejar a distancia y suma riesgo \\
Tema oscuro por defecto & La central está en penumbra durante toda la jornada \\
Detener con pulsación sostenida & Es la orden que más conviene que llegue y la que menos conviene mandar por un toque accidental \\
El repositorio simulado habla JSON & Así el contrato se prueba en la app y no queda solo en el documento \\
Inyección manual con un \texttt{AppContainer} & Para el tamaño de la app alcanza; Hilt sumaría configuración sin aportar \\
Sin imágenes en el historial & Privacidad de las personas que aparecen en las rondas \\
\bottomrule
\end{tabularx}
\end{table}

\subsection{Limitaciones}

\begin{itemize}
  \item No hay robot ni gateway: todo el comportamiento lo genera el simulador, con reglas simples y deterministas.
  \item La detección de personas y de puertas abiertas se simula. En un robot real haría falta un modelo de visión propio.
  \item El plano es una simplificación de una planta baja; no corresponde a la planta real de la sede.
  \item Hay un solo operador. No contemplamos dos guardias mandando órdenes al mismo robot.
  \item \pendiente{completar con lo que surja durante la implementación}
\end{itemize}
